\section{発展問題}\label{節:総合発展}

\begin{enumerate}[label=\textbf{\thechapter.\arabic*},leftmargin=*,start=13]
 \item\label{問:総合十三積分方程式}
 $[0,1]$で連続な関数$f$が,すべての$x\in[0,1]$について
 \[
  f(x)=1+\int_0^x(x-t)f(t)\dd t
 \]
 を満たすとする.\,$f$を求め,求めた関数が確かにこの等式を満たすことを確かめよ.
 微分方程式の一般的な解法は用いず,積の微分と基本定理から求めてよい.
 \item\label{問:総合十四二つの制約}
 $[0,1]$で連続な実数値関数$f$が
 \[
  \int_0^1f(x)\dd x=1,\qquad
  \int_0^1xf(x)\dd x=0
 \]
 を満たすとする.\,$\int_0^1f(x)^2\dd x$の最小値と,それを与える関数を求めよ.
 \item\label{問:総合十五媒介曲線}
 $0\le t\le\pi$で
 \[
  x=t-\frac12\sin2t,\qquad y=\sin^2t
 \]
 と表される曲線を考える.この曲線と$x$軸で囲まれる領域の面積,曲線の長さ,領域を$x$軸のまわりに一回転してできる立体の体積を求めよ.
 面積と体積の積分で,同じ部分を重複して数えていないことも確かめること.
 \item\label{問:総合十六近似と条件収束}
 次の2つの問いに,原始関数の初等的な表示を求めずに答えよ.
 \begin{enumerate}[label=(\arabic*),leftmargin=*]
  \item 有限テイラー公式を使い,\,$A=\int_0^1\sin(x^2)\dd x$を幅$\frac1{10000}$以下の有理数の区間に挟め.
  \item 広義積分$\int_1^{\infty}\sin(x^2)\dd x$が条件収束することを示せ.
 \end{enumerate}
 \item\label{問:総合十七ウォリス}
 非負整数$n$に対して
 \[
  C_n=\int_0^1(1-x^2)^n\dd x
 \]
 とする.\,$C_0$と$C_n$の漸化式を求め,\,$C_n$を有限の積で表せ.
 次にウォリスの積分との関係を使って,
 \[
  \lim_{n\to\infty}\sqrt n\,C_n,
  \qquad
  \lim_{n\to\infty}n^{\frac32}(C_n-C_{n+1})
 \]
 を求めよ.
 \item\label{問:総合十八相補と半角}
 $n$を正の整数とする.次の広義積分を,\,$n$を含む有限の積で表せ.
 \[
  J_n=\int_{-\pi}^{\pi}
  \frac{\sin^{2n}x}{(1+\e^x)\sqrt{1+\cos x}}\dd x.
 \]
 最初に両端での収束を確かめること.対称性,半角公式,置換積分,部分積分を組み合わせれば,ベータ関数やガンマ関数の公式を使わずに求められる.
\end{enumerate}

\clearpage

